\subsection{主要定理的叙述}

设 X 是一个集合。若 \(m \leqslant n\)，我们将用 \(\pi_{m,n}\) 表示从 \(X^{n}\) 到 \(X^{m}\) 中的、使得 \(\pi_{m,n}(\boldsymbol{x}) = (x_{1}, \ldots, x_{m})\)（对一切 \(\boldsymbol{x} = (x_{1}, \ldots, x_{n}) \in \mathrm{X}^{n}\)）的映射。

设 A 是 C 上的一个幺 Banach 代数。对每个整数 \(n \geqslant 1\) 及每个 \(a \in A^{n}\)，我们用 \(\mathrm{Sp}^{n}(\boldsymbol{a})\) 表示联合谱 \(\mathrm{Sp}_{\mathrm{A}}^{\{1,\ldots,n\}}(\boldsymbol{a})\) (I, p.~42, 定义 2)。它是 \(C^{n}\) 的一个紧子集。对每个满足 \(1 \leqslant m \leqslant n\) 的整数 m，有 \(\pi_{m,n}(\mathrm{Sp}^{n}(\boldsymbol{a})) = \mathrm{Sp}^{m}(\pi_{m,n}(\boldsymbol{a}))\) (I, p.~41, 第 6 小节)。连续线性映射

\[
\pi_ {m, n} ^ {*} \colon \mathcal {O} (\mathrm{Sp} ^ {m} (\pi_ {m, n} (\boldsymbol {a})); \mathrm{A}) \longrightarrow \mathcal {O} (\mathrm{Sp} ^ {n} (\boldsymbol {a}); \mathrm{A})
\]

是幺代数的一个态射。

设 A 是 C 上的一个交换幺 Banach 代数。设 \(n \geqslant 1\) 是一个整数。我们把这样的资料称为 A 上 n 个变量的全纯函数演算：对每个 \(a \in A^{n}\)，给定一个映射

\[
\Theta_ {\boldsymbol {a}} \colon \mathcal {O} (\mathrm{Sp} ^ {n} (\boldsymbol {a}); \mathrm{A}) \longrightarrow \mathrm{A}
\]

满足下列条件：

(CF1) 对每个 \(a \in A^{n}\)，映射 \(\Theta_{a}\) 是幺代数的一个连续态射。

(CF2) 若 \(\boldsymbol{a}=(a_{1},\ldots,a_{n})\)，且 \(z_{1},\ldots,z_{n}\) 表示 \(\mathrm{Sp}^{n}(\boldsymbol{a})\) 的邻域中 \(C^{n}\) 的坐标函数的芽，则有

\[
\Theta_ {\boldsymbol {a}} (z _ {1}) = a _ {1}, \dots , \Theta_ {\boldsymbol {a}} (z _ {n}) = a _ {n}.
\]

注. --- 若代数 A 的根为零，则可以省去 (CF1) 中的连续性条件（参见 I, p.~40, 命题 9）。

我们把这样的资料称为 A 上的全纯函数演算：对每个整数 \(n \geqslant 1\)，给定 A 上 n 个变量的一个全纯函数演算，且满足：

(CF3) 无论 m 与 n 是怎样的整数，只要 \(1 \leqslant m \leqslant n\)，也无论 \(a \in A^{n}\) 与 \(f \in \mathcal{O}(\mathrm{Sp}^{m}(\pi_{m,n}(a); \mathrm{A})\) 是怎样的，总有

\[
\Theta_ {\boldsymbol {a}} (\pi_ {m, n} ^ {*} (f)) = \Theta_ {\pi_ {m, n} (\boldsymbol {a})} (f).
\]

本节的目的在于证明下述定理：

定理 1. --- 设 A 是一个复交换幺 Banach 代数。则存在唯一的一个 A 上的全纯函数演算。

这一定理的证明将占据第 3 至第 7 小节。

